%% kapa: exam=1 round=29 year=2018 period=1 form=A subject=economics subject_ko=경제학원론 seq=32 number=72
%% source: Q-Net 감정평가사 기출문제 (cst00309) artlSeq=5202808 file=extracted/1008/5202808/2018년 제29회 감정평가사 1차시험 기출문제/2018년 제29회 감정평가사 1차시험 기출문제/1교시 A형.pdf
%% conversion: kordoc 4.24.51 -> tools/gen1.py (auto); review=proofread-sonnet; flags=none
\begin{question}{72}
甲국 통화당국의 손실함수와 필립스곡선이 다음과 같다. 인플레이션율에 대한 민간의 기대가 형성되었다. 이후, 통화당국이 손실을 최소화하기 위한 목표 인플레이션율은? (단, $\pi ,\pi^{e},u,u_{n}$ 은 각각 인플레이션율, 민간의 기대인플레이션율, 실업률, 자연실업률이고, 단위는 \%이다.)
\begin{framedtext}
통화당국의 손실함수: $L(\pi ,u)=u+\frac{1}{2}\pi^{2}$ \par
필립스곡선: $\pi =\pi^{e}-\frac{1}{2}(u-u_{n})$ \par
\end{framedtext}
\begin{choices}
  \item 0 \%
  \item 1 \%
  \item 2 \%
  \item 3 \%
  \item 4 \%
\end{choices}
\end{question}
\answer{3}
